%--------------------------------------
% Caso de Uso: Eliminar Reservación
%--------------------------------------

\begin{UseCase}{CU17}{Eliminar Reservación}{
	Permitir que un usuario autorizado cancele o elimine una reservación existente en el sistema del Hotel para Perros.
}
	\UCitem{Versión}{\color{Gray}1.0}
	\UCitem{Autor}{\color{Gray}Castillo Aldaba Jared}
	\UCitem{Supervisa}{\color{Gray}Ugalde Tellez Aaron}
	\UCitem{Actor}{\hyperlink{Administrador}{Administrador}}
	\UCitem{Propósito}{Cancelar una reservación eliminándola del sistema o cambiando su estado a cancelada.}
	\UCitem{Entradas}{Reservación seleccionada, motivo de cancelación opcional.}
	\UCitem{Origen}{Pantalla de detalle de reservación}
	\UCitem{Salidas}{Reservación cancelada y habitaciones liberadas.}
	\UCitem{Destino}{Pantalla de reservaciones y base de datos}
	\UCitem{Precondiciones}{
		\begin{itemize}
			\item El usuario debe haber iniciado sesión con permisos adecuados.
			\item La reservación debe existir en el sistema.
			\item La reservación debe encontrarse en un estado cancelable.
		\end{itemize}
	}
	\UCitem{Postcondiciones}{
		\begin{itemize}
			\item La reservación queda marcada como cancelada.
			\item Las habitaciones asignadas quedan disponibles.
			\item Los servicios no realizados se liberan.
			\item El calendario se actualiza correctamente.
		\end{itemize}
	}
	\UCitem{Errores}{
		Reservación no encontrada, permisos insuficientes, estado no cancelable, error al actualizar la información.
	}
	\UCitem{Tipo}{Caso de uso secundario}
	\UCitem{Observaciones}{
		\begin{itemize}
			\item Las reservaciones en curso no pueden cancelarse.
			\item Puede aplicarse una política de reembolso según la fecha de cancelación.
			\item Se registra auditoría de la operación realizada.
		\end{itemize}
	}
\end{UseCase}

%--------------------------------------
\begin{UCtrayectoria}
	\UCpaso[\UCactor] Selecciona una reservación desde la lista o el calendario.
	\UCpaso[\UCactor] Accede a la pantalla de detalle de la reservación.
	\UCpaso[\UCactor] Selecciona la opción para cancelar la reservación.
	\UCpaso El sistema solicita confirmación y permite ingresar un motivo opcional.
	\UCpaso[\UCactor] Confirma la cancelación.
	\UCpaso El sistema verifica los permisos del usuario.
	\UCpaso El sistema valida que la reservación pueda cancelarse.
	\UCpaso El sistema inicia la operación de actualización.
	\UCpaso El sistema marca la reservación como cancelada.
	\UCpaso El sistema libera las habitaciones asociadas.
	\UCpaso El sistema cancela los servicios futuros.
	\UCpaso El sistema registra la operación realizada.
	\UCpaso El sistema confirma la operación.
	\UCpaso El sistema muestra un mensaje indicando que la reservación fue cancelada correctamente.
	\UCpaso El sistema actualiza la lista y el calendario de reservaciones.
\end{UCtrayectoria}

%--------------------------------------
\begin{UCtrayectoriaA}{A}{Reservación no encontrada}
	\UCpaso El sistema muestra un mensaje indicando que la reservación no existe.
	\UCpaso El sistema regresa a la lista de reservaciones.
\end{UCtrayectoriaA}

%--------------------------------------
\begin{UCtrayectoriaA}{B}{Permisos insuficientes}
	\UCpaso El sistema informa que el usuario no tiene permisos para cancelar reservaciones.
	\UCpaso El caso de uso termina sin realizar cambios.
\end{UCtrayectoriaA}

%--------------------------------------
\begin{UCtrayectoriaA}{C}{Reservación en estado no cancelable}
	\UCpaso El sistema informa que la reservación no puede cancelarse en su estado actual.
	\UCpaso El caso de uso termina sin realizar cambios.
\end{UCtrayectoriaA}

%--------------------------------------
\subsection{Puntos de extensión}
\UCExtenssionPoint{
	El usuario desea aplicar una política de reembolso por cancelación.
}{
	Del paso 12.
}{
	Procesar Reembolso por Cancelación.
}

%--------------------------------------
% Fin del Caso de Uso
